%%%%%%%%%%%%%%%%%%%%%%%%%%%%%%%%%%%%%%%%%%%%%%%%%%%%%%%%%%%%%%%%%%%%%%%%%%%%
% AGUJournalTemplate.tex: this template file is for articles formatted with LaTeX
%
% This file includes commands and instructions
% given in the order necessary to produce a final output that will
% satisfy AGU requirements, including customized APA reference formatting.
%
% You may copy this file and give it your
% article name, and enter your text.
%
% guidelines and troubleshooting are here:

%% To submit your paper:
\documentclass{agujournal2019}

\usepackage{url} %this package should fix any errors with URLs in refs.
\usepackage{lineno}
\usepackage[inline]{trackchanges} %for better track changes. finalnew option will compile document with changes incorporated.
\usepackage{soul}
\usepackage{amsmath}
\usepackage{amssymb}


\linenumbers

\draftfalse

\journalname{JGR: Earth Surface, Confidential manuscript under review}


\begin{document}

\title{Constraining the Distribution of 3D Fractal Structures in Mud Flocs}
\authors{
Brayden Noh\affil{1,2},
Justin A.~Nghiem\affil{1},
Kimberly Litwin Miller\affil{1},
Dongchen Wang\affil{1},
Michael P.~Lamb\affil{1}
}

\affiliation{1}{Division of Geological and Planetary Sciences, California Institute of Technology, Pasadena, CA 91125, USA}
\affiliation{2}{Department of Earth and Planetary Sciences, Harvard University, Cambridge, MA 02138, USA}

\correspondingauthor{Brayden Noh}{braydennoh@fas.harvard.edu}


\begin{keypoints}
\item Aggregating floc data overestimates 3D fractal dimension due to structural diversity.
\item Stratifying by 2D fractal dimension reveals increasing $n_f$ with shear and decreasing $n_f$ with organic matter.
\item Structure-resolved $n_f$ estimates yield settling velocities 2--10$\times$ lower than conventional aggregated estimates, implying longer transport distances.
\end{keypoints}




\begin{abstract}
Mud builds coastal landscapes and hosts pollutants and particulate carbon, yet predicting its transport is difficult because mud aggregates into flocs, which control settling rates in riverine, estuarine, and marine environments. Flocs are typically characterized using fractal theory with a constant fractal dimension ($n_f = 2.0$--$2.2$), inferred from aggregated floc size--settling velocity data despite their complex fractal structures. We conducted experiments on freshwater flocs under varied shear stress and organic matter concentrations to characterize this diversity. The fractal dimension $n_f$ was calculated from coupled measurements of floc size and settling rate using particle image tracking. Rather than the conventional approach of aggregating all settling data to infer a single value of $n_f$, we grouped measurements by the two-dimensional fractal dimension derived from box counting, thereby partially resolving structural variability. This stratified approach reveals that $n_f$ increases with shear velocity (from 1.77 to 2.12) and decreases with organic matter content (from 1.84 to 1.63), trends that are obscured when data are aggregated. Across our conditions, aggregation biases inferred $n_f$ by $\sim$5--10\%, which propagates nonlinearly into settling velocities and predicted transport distances. The relationship between 2D and 3D fractal dimensions is best described by a power-law conversion ($n_f = 0.81\,(n_f^{(2\mathrm{D})})^{1.81}$; \citeA{wang2022fractal}), consistent with the sub-Euclidean space filling of porous aggregates. Although conventional bulk fits to our data yield $n_f$ values near the canonical range, we show that these estimates are biased due to correlations among settling velocity, floc size, and fractal dimension, analogous to Simpson's Paradox. Our approach reduces this bias, yields lower $n_f$, and implies that mud flocs are more porous, settle more slowly, and travel farther than previous models predict. This finding suggests that freshwater mud retention in floodplains and deltas is less efficient than current floc models assume---at least under the freshwater conditions analyzed here---with implications for sediment budgets, contaminant dispersal, and coastal landscape evolution.


\end{abstract}

\section*{Plain Language Summary}
Mud particles in rivers tend to clump together into clusters called flocs. How quickly these flocs sink determines where mud is deposited, shaping the growth of river deltas, floodplains, and coastlines. Predicting this settling is challenging because flocs vary widely in structure: some are dense and compact, while others are porous and loosely bound. Most studies assume that all flocs have a single, uniform structure, which oversimplifies their behavior. In our experiments, we generated flocs in freshwater under controlled conditions with different flow strengths and amounts of organic matter. Using high-resolution imaging, we measured both their structure and settling speed. We found that faster flows tend to produce more compact flocs that settle more quickly, while organic matter encourages the formation of more open structures that settle more slowly. By grouping flocs according to their measured structural complexity, we revealed trends that are hidden when data are grouped together. This new approach reduces bias in estimating settling rates and improves predictions of how mud moves through rivers and into estuaries and oceans, where salinity may further modify physical clumping beyond the freshwater trends isolated here.

\section{Introduction}

Fine-grained sediment, or mud, defined by particles smaller than 62.5~\textmu m, dominates fluvial sediment loads and represents the primary material deposited in lowland and deltaic environments \cite{healy2002muddy, kranenburg1994fractal, lamb2020mud, nghiem2026flocculated}. The dynamics of mud transport and deposition therefore exert first-order control on the long-term evolution of deltas, floodplains, and estuaries \cite{nghiem2022mechanistic, winterwerp2002flocculation}, as well as on water quality, contaminant fluxes, and the success of restoration interventions \cite{droppo1997freshwater, esposito2017efficient, kemp2016enhancing, liss1996floc}. Mud also plays a central role in the global carbon cycle by providing reactive mineral surfaces for organic carbon adsorption and long-term preservation \cite{bianchi2024anthropogenic, blair2012fate}.

The geomorphic and geochemical importance of mud necessitates a mechanistic framework for particle transport and settling to enable predictive modeling. Settling velocity is a key parameter in sediment-transport equations and is commonly estimated using Stokes' law, which gives the terminal settling velocity ($w_s$) of smooth, impermeable spheres at low Reynolds numbers \cite{bird2006transport, stokes1851effect}:
\begin{equation}
w_s = \frac{g\,(\rho_p - \rho_w)\,d^2}{18\,\rho_w\,\nu}
    \label{eq:stokes}
\end{equation}
where $g = 9.81\ \mathrm{m\,s^{-2}}$ is gravitational acceleration, $\rho_p\ (\mathrm{kg\,m^{-3}})$ and $\rho_w\ (\mathrm{kg\,m^{-3}})$ are the densities of the primary particle and water, respectively, $d\ (\mathrm{m})$ is particle diameter, and $\nu\ (\mathrm{m^2\,s^{-1}})$ is the kinematic viscosity of the water.

However, when applied to fine-grained sediment, Stokes' law yields predictions at odds with field observations. Consider a single 1~\textmu m primary mineral grain (e.g., quartz or clay, $\rho_p \approx 2650\ \mathrm{kg\,m^{-3}}$) in a $5$\,m-deep river flowing at $1\,\mathrm{m\,s}^{-1}$. Using Stokes' law produces a settling velocity of approximately $9\times 10^{-7}\,\mathrm{m\,s}^{-1}$, implying a deposition time of about 64 days and a transport distance exceeding 5,500\,km downstream. This hypothetical scenario conflicts sharply with field observations of mud deposition \cite{lamb2020mud, mason2019differential, nicholas1996significance}. Field studies document that lowland rivers commonly build levees through overbank flows, with many levees composed primarily of mud and clay contents of approximately $50\%$ \cite{aalto2008spatial}. Furthermore, deposition rates along levees typically decrease exponentially with distance from the channel, indicating transport lengths on the order of tens to hundreds of meters rather than thousands of kilometers \cite{hajek2012simplified}.

This discrepancy indicates that mud settles far more rapidly than predicted for isolated particles, implying the presence of additional physical processes. The deviation arises because suspended mud particles flocculate, forming larger aggregates that are composed of individual primary particles (representing the constituent mineral grains) and that settle much faster than their isolated constituents \cite{droppo1997freshwater, johnson1996settling, kranenburg1994fractal}. These flocculated aggregates, or flocs, are highly porous rather than dense, impermeable spheres, so their effective density decreases with increasing size \cite{kranenburg1994fractal, Meakin1992Aggregation}. Because flocs exhibit approximately self-similar fractal organization \cite{family2012kinetics, jullien1987aggregation}, their hierarchical structure can be described by a standard fractal mass--size scaling \cite{Meakin1992Aggregation}:
\begin{equation}
N \sim \left( \frac{d_f}{d_p} \right)^{n_f},
    \label{eq:fractaltheory}
\end{equation}
where $N$ is the number of primary particles (dimensionless), $d_f$ is the floc diameter (m), $d_p$ is the primary particle diameter (m), and $n_f$ is the 3D fractal dimension (dimensionless).

Building on Eq.~\eqref{eq:fractaltheory}, the fraction of the floc volume actually occupied by solid grains is the ratio of solid volume to the enclosing floc volume. Because the solid volume scales with $N$ times the monomer volume while the enclosure scales with $d_f^{3}$, the solid volume fraction $\phi_s$ decreases with size approximately as $(d_f/d_p)^{\,n_f-3}$. For a porous aggregate whose pores are filled by ambient water, the floc's submerged specific gravity equals this solid fraction times the primary particle's submerged specific gravity. Denoting submerged specific gravity by $R=(\rho_p-\rho_w)/\rho_w$, this yields:
\begin{equation}
R_f = R_p \left( \frac{d_f}{d_p} \right)^{n_f - 3},
    \label{eq:eq3}
\end{equation}
where $R_f$ and $R_p$ are the submerged specific gravities of the floc and primary particles, respectively (both dimensionless), $\rho_p$ is primary particle density (kg\,m$^{-3}$), and $\rho_w$ is water density (kg\,m$^{-3}$). Since a dense solid sphere has a volume that scales with the cube of its diameter, its fractal dimension is $n_f=3$, and Eq.~\eqref{eq:eq3} reduces to $R_f=R_p$. Naturally formed flocs are more porous with $1 \le n_f < 3$, yielding $R_f<R_p$. Accordingly, floc effective (excess) density decreases with floc size. For a fractal aggregate, the solids mass scales as $M_f \propto d_f^{n_f}$ whereas the envelope volume scales as $V_f \propto d_f^{3}$, implying the excess density, $\Delta \rho_f$, scales as $\Delta \rho_f \propto M_f / V_f \propto d_f^{n_f-3}$ \cite{kranenburg1994fractal}.

Substituting Eq.~\eqref{eq:eq3} into Eq.~\eqref{eq:stokes} gives an explicit expression for the settling velocity of a fractal floc \cite{strom2011explicit, winterwerp1998simple}:
\begin{equation}
    w_s = \frac{g\,R_p}{b_1\,\nu\,d_p^{\,n_f - 3}}\; d_f^{\,n_f - 1},
    \label{eq:floc_settling}
\end{equation}
where $w_s$ is the settling velocity (m\,s$^{-1}$) and $b_1$ is a dimensionless shape factor (with $b_1=18$ for smooth spheres).

In practice, applying Eq.~\eqref{eq:floc_settling} requires an appropriate fractal dimension $n_f$, which captures how floc structure modifies settling velocity. Experimental and field studies often assume $n_f \approx 2.0$--$2.2$ from empirical $w_s$--$d_f$ relationships, values reported primarily for saline estuarine/coastal flocs where electrochemical interactions promote more compact aggregates; freshwater flocs more often show lower $n_f$ ($\lesssim 2$), with overlap depending on turbulence and composition \cite{winterwerp1998simple,winterwerp2006heuristic}. With all other parameters held fixed, Eq.~\eqref{eq:floc_settling} gives $w_s \propto d_f^{n_f-1}$ \cite{johnson1996settling, li1989fractal}. Measurements of $w_s$ and $d_f$ across a floc population therefore allow $n_f$ to be inferred from the log--log slope of the $w_s$--$d_f$ relation \cite{kranenburg1994fractal, strom2011explicit, winterwerp1998simple, winterwerp2002flocculation, winterwerp2004introduction, winterwerp2006heuristic}. In general, this inference assumes that floc permeability is independent of floc size. If permeability instead varies systematically with $d_f$, the observed scaling relation may reflect both fractal structure and permeability effects \cite{strom2011explicit, nghiem2026flocculated}.

Using a constant population-wide value for the fractal dimension is convenient but restrictive. \citeA{khelifa2006models} showed that treating $n_f$ as constant fails to reproduce the observed spread in settling velocities and argued that $n_f$ can decrease with floc size, motivating models that prescribe a function $n_f(d_f)$, often as a power law, to represent within-population variability. In practice, however, most experimental and image-based field studies measure only floc size $d_f$ and settling velocity $w_s$ \cite{smith2015image, nghiem2024testing}, so the fractal dimension is not independently observed and functions for $n_f$ cannot be directly tested. Moreover, models for $n_f(d_f)$ implicitly assume a unique fractal dimension for a given floc size \cite{khelifa2006models}, whereas natural floc populations can contain multiple flocs with comparable $d_f$ but distinct internal structure, and therefore distinct $n_f$ \cite{maggi2007variable}. As a result, standard log--log regressions of $w_s$ on $d_f$ implicitly treat structurally heterogeneous populations as homogeneous, and data aggregation of such subpopulations can misattribute between-group differences to size scaling itself, producing biased estimates of $n_f$.


To illustrate this effect, consider a population comprising three subgroups of flocs, each defined by a distinct fractal dimension ($n_f = 1.5$, $2.0$, and $2.5$) but formed from primary particles of the same diameter ($d_p = 10^{-5}$~m). According to Eq.~\eqref{eq:floc_settling}, the slope of the log--log relationship between settling velocity and floc diameter for each subgroup is $n_f - 1$. When analyzed separately, these subgroups yield distinct log--log slopes corresponding to $n_f = 1.5$, $2.0$, and $2.5$, respectively (Fig.~\ref{fig:1}). However, if these floc subgroups are combined into a single group and floc diameter ($d_f$) is correlated with the 3D fractal dimension ($n_f$), then fitting a single power-law relationship to all points can yield a substantially skewed and misleading exponent. In Fig.~\ref{fig:1}, the resulting log--log slope is $2$, implying an effective $n_f = 3$, which corresponds to the settling behavior of smooth, impermeable spheres, not flocs. This example parallels Simpson's Paradox \cite{simpson1951interpretation}, a statistical phenomenon in which relationships observed within individual subgroups are reversed, masked, or distorted upon data aggregation. In the context of floc settling, data aggregation of structurally distinct subpopulations can produce a global $w_s$--$d_f$ slope that implies an apparent $n_f$ inconsistent with the $n_f$ of any subgroup, so the inferred ``representative'' fractal dimension is a regression artifact rather than a physical descriptor.

\begin{figure}[htbp]
    \centering
    \includegraphics[width=0.8\textwidth]{figures/1.pdf}
    \caption{
    Synthetic data illustrating Simpson's Paradox in floc settling. Three floc subgroups, each with a specific 3D fractal dimension ($n_f = 1.5,\ 2.0,\ \text{and}\ 2.5$), are shown. By Eq.~\eqref{eq:floc_settling}, the slope of the log--log relationship between settling velocity ($w_s$) and floc diameter ($d_f$) is $n_f - 1$. Fitting each subgroup individually recovers slopes of 0.5, 1.0, and 1.5, as expected. In contrast, data aggregation of all points produces a single power-law fit (red line) with a misleading slope of 2, equivalent to $n_f = 3$, characteristic of solid spheres and not flocs. All trend lines intersect at the primary particle diameter ($d_f = d_p = 10^{-5}\,\mathrm{m}$), where settling velocity is independent of fractal dimension.
    }
    \label{fig:1}
\end{figure}

To estimate fractal dimension independently of settling velocity and floc diameter, studies have employed image-based techniques that analyze in situ 2D projections of flocs \cite{jarvis2005measuring}. These approaches typically calculate a two-dimensional (2D) fractal dimension from high-resolution images, then infer the corresponding three-dimensional (3D) structure via empirical or simulation-based relationships. A common method relates the perimeter--area scaling of projected flocs to a 2D fractal dimension, using computationally generated aggregates to derive a mapping between 2D and 3D fractal dimension \cite{lee2004prediction, maggi2004method, tang2015reconstructing}. However, the relationship between 2D and 3D fractal dimension is derived from synthetic aggregates, which may not replicate the aggregation dynamics or structural variability of natural flocs. Projection-based estimates therefore risk systematic bias and, more importantly, do not recover the 3D fractal dimension that directly controls settling behavior.

Although flocculation has been widely characterized, how within-population structural heterogeneity affects settling remains uncertain, particularly under natural particle-size distributions and aggregation pathways. The central difficulty is that $w_s - d_f$ regressions alone do not condition on structure, so data aggregation of distinct subpopulations can produce an apparent relation and inferred $n_f$ that is biased by Simpson-type effects. To address this bias, we use images from experiments to stratify the floc population by the two-dimensional (2D) fractal dimension from box-counting, and within each subgroup, regress settling velocity $w_s$ against floc diameter $d_f$ to estimate the hydrodynamically relevant three-dimensional (3D) fractal dimension $n_f$. This conditioning step provides a principled way to separate structural heterogeneity from size scaling without prescribing a population-level closure such as $n_f(d_f)$: that is, $n_f$ is inferred hydrodynamically within approximately homogeneous structural classes. The assumption is that flocs with similar 2D fractal dimension have similar 3D fractal dimension, which we validate herein.

We evaluate this framework in controlled laboratory experiments designed to simulate flocs in rivers with systematic variation in shear velocity ($u_*$) and organic matter concentration. High-resolution in situ imaging enabled quantification of floc morphology and settling trajectories. Specifically, we pursued three objectives: (1) demonstrate the existence of structurally distinct subgroups within the floc population, each characterized by a unique fractal dimension; (2) evaluate the relationship between 2D and 3D fractal dimension, and assess this relationship against theoretical predictions; and (3) estimate the effective primary particle size ($d_p$) and quantify how the hydrodynamic shape factor ($b_1$) varies as a function of $n_f$. All experiments were conducted in freshwater (salinity $\approx 0$), so the results isolate the roles of shear and organic binding and are not intended as a direct constraint for saline estuaries where salt-driven cohesion can further modify flocculation.

\section{Methodology}

\subsection{Experimental Setup}

We performed flocculation experiments in a custom mixing tank designed to provide controlled conditions (Fig.~\ref{fig:2}a). The tank comprises a cylindrical PVC pipe, 20~cm in diameter and 60~cm tall, filled with water to a depth of 40~cm. A motor-driven rotating lid, mounted above the open top and submerged just below the free surface, sets the desired shear. Shear velocity, $u_*$, was empirically calibrated for each run by converting motor revolutions per minute to $u_*$ using a relationship derived from previous sediment-entrainment trials in \citeA{douglas2025mud}. We define $u_* = (\tau/\rho_w)^{1/2}$, where $\tau$ is a characteristic turbulent fluid shear stress and $\rho_w$ is the density of water, and we treat $u_*$ as an effective, run-averaged measure of turbulent shear intensity because the rotating-lid-driven flow field is spatially heterogeneous.

During the mixing phase, the flow was fully turbulent. This is quantified by the friction Reynolds number $Re_{\tau} = u_* h / \nu$, where $h$ is the water depth and $\nu$ is the kinematic viscosity of freshwater. For $h = 0.40$~m and $\nu \approx 1.0 \times 10^{-6}\ \mathrm{m^2\,s^{-1}}$ at room temperature, the imposed shear velocities $u_* = 0.02$--$0.04$~m~s$^{-1}$ correspond to $Re_{\tau} \approx (8$--$16)\times10^{3}$, well within the fully turbulent regime for wall-bounded shear flows \cite{Pope2000,Schlichting2000}. Across the flocculation experiments, we imposed $u_* = 0.02$--$0.04$~m~s$^{-1}$ (equivalent to $\tau \approx 0.4$--1.6~Pa), spanning low-to-moderate shear conditions typical of suspension-dominated river channels, sufficient for sand entrainment but below those required to mobilize gravels \cite{julien2018river}.

Suspended flocs were imaged through a 3~mm-wide flow-through slit built into the tank wall at the observation level. A DSLR fitted with a $5\times$ microscope objective was aligned with the slit to record flocs as they crossed the narrow optical path, enhancing image clarity and increasing the number of focused particles captured by the camera, following methods similar to \citeA{osborn2021flocarazi}.

\begin{figure}[htbp]
    \centering
    \includegraphics[width=1\textwidth]{figures/2.pdf}
    \caption{
    (a) Schematic of the custom flocculation apparatus. A cylindrical PVC tank (20~cm diameter, 60~cm height) is filled to 40~cm depth and stirred with a motor-driven lid just below the surface. Shear velocity is calibrated for each run. An observation slit (3~mm wide) enables \textit{in situ} imaging of suspended flocs with a DSLR camera and 5$\times$ microscope objective aligned for high-resolution focus. Key components and dimensions are annotated.
    (b) Grain-size distributions of the dispersed input sediment measured by Mastersizer from \citeA{nghiem2025flocculation}. Curves show pure silica sand (SiO$_2$), montmorillonite clay, and their 50:50 volume-weighted mixture.}
    \label{fig:2}
\end{figure}

\subsection{Experimental Runs}

We conducted eight experimental runs in total. The first two were calibration runs using laboratory spheres and silica grains, which do not flocculate, to verify the accuracy of our settling-velocity measurements and provide a baseline for comparing non-flocculating particles to flocs under identical experimental conditions. The remaining six runs systematically varied shear velocity ($u_*$) and organic matter (OM) content, given that previous studies have shown turbulent shear limits floc size by promoting breakup, while organic matter enhances physical aggregation by binding particles \cite{nghiem2022mechanistic}. By independently controlling these variables, we aimed to isolate their respective roles in floc formation and settling.

For each experiment, the tank was filled with water and the motor was started to establish the target shear. Particles and organic matter were then added gradually, and the tank was mixed for 30~min to allow flocs to form under sustained shear. The motor was subsequently stopped to initiate settling, and video recording began immediately and continued for the first 30~min. All measurements of floc size, settling velocity, and fractal dimension reported in this study were collected during this settling phase, not during mixing. Each floc experiment used 10~g of solids: 5~g silica (SiO$_2$) and 5~g montmorillonite clay. Organic content was varied in three runs by adding guar gum, a plant-derived polysaccharide used as a proxy for natural organic polymers, at 1, 2, or 3~wt\% relative to the total mass of solids, following the organic-loading framework of \citeA{zeichner2021early}. This range was selected to (i) remain within the lower end of organic fractions reported for natural river suspended sediment (on the order of 1--6\% and up to \(\sim\)20\% in compiled field datasets), and (ii) avoid polymer-dominated regimes while still producing a resolvable change in flocculation dynamics. In \citeA{zeichner2021early}, polymer-to-clay ratios were chosen explicitly to span modern river organic-matter concentrations, and their experiments show that even 1--2~wt\% guar gum produces a strong flocculation response. Guar gum is a controlled proxy for natural extracellular polymeric substances and does not capture the full chemical complexity of field organic matter. Full experimental conditions, including organic-matter fraction and imposed $u_*$, are listed in Table~\ref{tab:exp_conditions}.

\begin{table}[htbp]
\caption{Summary of experimental conditions.\label{tab:exp_conditions}}
\centering
\begin{tabular}{lccc}
\hline
Exp Num & Experiment Type & OM (wt\%) & Shear velocity (m/s) \\
\hline
1 & Lab Sphere     & NA    & NA    \\
2 & Silica         & NA    & NA    \\
3 & Floc Shear 1   & 2     & 0.02  \\
4 & Floc Shear 2   & 2     & 0.03  \\
5 & Floc Shear 3   & 2     & 0.04  \\
6 & Floc OM 1     & 1     & 0.04  \\
7 & Floc OM 2     & 2     & 0.04  \\
8 & Floc OM 3     & 3     & 0.04  \\
\hline
\multicolumn{4}{l}{NA: Not applicable.}
\end{tabular}
\end{table}

Grain-size distributions for the silica and montmorillonite were measured using a Malvern Mastersizer 3000 (Fig.~\ref{fig:2}b), following the procedure of \citeA{nghiem2025flocculation}. In brief, samples were dispersed in deionized water, sonicated to break up aggregates, and analyzed by laser diffraction to obtain particle-size distributions. Together, these two minerals span nearly three orders of magnitude in diameter, yielding a volume-weighted size distribution from $2 \times 10^{-7}$ to $2 \times 10^{-4}$~m. For the mixed-sediment runs, a composite distribution was calculated as the 50:50 volume-weighted average of the two minerals.

\subsection{Image Processing Methods for Floc Settling Velocity and Diameter}

The floc tracking analysis began with processing every frame of the recorded video sequence, acquired at a constant frame interval $\Delta t = 1/30~\mathrm{s}$, as verified from the video metadata. Each frame was converted to an eight-bit grayscale image using OpenCV, and a time-averaged background image (computed by averaging a block of 1000 consecutive frames) was subtracted from each image to isolate transient features such as moving flocs (Fig.~\ref{fig:3}a).

\begin{figure}[htbp]
    \centering
    \includegraphics[width=1\textwidth]{figures/3.pdf}
    \caption{
    Image processing workflow used for floc analysis. (a) Example of a background-subtracted image, highlighting transient objects such as flocs. (b) Detected and contoured particles passing focus and sharpness thresholds (outlined in red) with raw velocity vectors overlaid. (c) Local background flow field calculated by PIV from tracer particles. (d) Example of a tracked floc with both its raw and corrected (background-subtracted) settling velocity vectors indicated.
    }
    \label{fig:3}
\end{figure}

Particle boundaries in each image were detected from intensity gradients using a global threshold of 200 applied to the normalized residual image. To retain only sharp and well-focused flocs, candidate objects were required to exceed a Laplacian variance threshold of 5 (measuring focus/sharpness) and to have an intensity standard deviation greater than 5 within their contour \cite{pertuz2013analysis}. These specific threshold values were optimized for our experimental conditions and image resolution, and may vary between different setups. Only particles meeting both criteria were accepted for further analysis and are highlighted in red in Fig.~\ref{fig:3}b. This objective per-frame filtering mitigates focus-related sampling bias by excluding blurred detections rather than allowing size-dependent blur to enter the floc statistics. For each accepted object, we recorded the planform area in image pixels $A_{\mathrm{px}}$ (later converted to an equivalent diameter after spatial calibration) and the centroid coordinates $(x_c,y_c)$ for subsequent velocity estimation.

To express diameters and velocities in physical units, we calibrated pixel size using a precision stage micrometer imaged under the same optical conditions as the experiments. Comparing the pixel spacing between micrometer markings with their true separations yielded a scale factor of \(0.45~\mu\mathrm{m\,px^{-1}}\). We applied this factor to all subsequent analyses to convert particle dimensions and inter-frame displacements from pixels to micrometers; velocities were then obtained from these displacements using the recorded frame interval.

Settling velocities were obtained by frame-to-frame tracking using a conservative nearest-neighbor association. For each detection in frame \(t\), we searched within a 50-pixel radius in frame \(t{+}1\) for candidate matches. We scored each candidate \(j\) using two equally weighted, normalized differences: the equivalent concave-hull diameter \(d_{\mathrm{ch}}\) (px) and the Laplacian-variance sharpness \(L\) (arb. units). The score was defined as \(S_j=\tfrac{1}{2}\big(|d_{\mathrm{ch},j}-d_{\mathrm{ch},t}|/d_{\mathrm{ch},t}+|L_j-L_t|/L_t\big)\). We selected the candidate with the smallest \(S_j\), provided that both relative differences were \(\le 0.10\); otherwise, the track was terminated to avoid spurious links from rapid shape or focus changes. The vertical velocity for that interval was initially computed as \(\Delta y\,s/\Delta t\), where \(\Delta y\) is the vertical centroid shift (px, positive downward), \(s\) is the spatial calibration (\(0.45~\mu\mathrm{m\,px^{-1}}\)), and \(\Delta t\) is the frame interval (s). Raw velocities computed in \(\mu\mathrm{m\,s^{-1}}\) were subsequently converted to \(\mathrm{m\,s^{-1}}\) to match the physical units required by Eq.~\eqref{eq:stokes} and Eq.~\eqref{eq:floc_settling}.

Frame-to-frame particle displacements yield raw two-dimensional velocity vectors that combine gravitational settling with background flow in the imaging plane. To isolate the settling signal, we first computed a local 2D flow field using OpenPIV \cite{liberzon2021openpiv} for each frame pair. Following \citeA{smith2015image}, the images were digitally filtered to retain only tracers smaller than $10~\mu\mathrm{m}$ that behave as passive flow followers (Fig.~\ref{fig:3}c). For each tracked floc, we then subtracted the corresponding local flow vector---bilinearly interpolated at the floc centroid---from the floc's raw trajectory vector. This vector differencing produces a background-corrected motion (Fig.~\ref{fig:3}d, gray arrow for raw, white arrow for corrected), from which we report the vertical component as the settling velocity \(w_s\). All settling velocities in this study refer to this corrected, background-subtracted quantity.

After velocity correction, floc diameters were derived from each particle's projected area. Two common image-based metrics are used to estimate diameter: the Feret diameter (maximum caliper length across the projection) and the equivalent circular diameter (ECD), defined as the diameter of a circle with the same projected area \cite{jarvis2005measuring, smith2015image}. The Feret diameter is sensitive to orientation and outliers, often overestimating size for irregular or elongated flocs. By contrast, the area-equivalent diameter---especially when computed from a concave hull---can better represent irregular shapes by de-emphasizing extreme protrusions, though it may underrepresent thin filaments. Accordingly, we adopt the concave-hull ECD (with $\alpha = 0.05$, the concave-hull radius parameter that controls the allowed boundary indentation), corresponding to the projected-area--equivalent diameter of \citeA{smith2015image}, while acknowledging that alternative definitions are reasonable \cite{jarvis2005measuring}. Particles were segmented with binary masks, boundaries were traced using a concave-hull algorithm, the enclosed area was measured, and the diameter of a circle of equal area was computed and used as \(d_f\) in all subsequent log--log analyses. Because $d_f$ is defined geometrically from the 2D projection, this step does not assume constant floc density; any density changes associated with size or organic content enter only through the measured settling velocities.

\subsection{Calibration and Settling Behavior Analysis}

We calibrated our setup with Experiments~1 and~2, which measured the settling of laboratory spheres (50~\textmu m) and natural silica grains (Fig.~\ref{fig:4}a). Settling velocities were fitted with Stokes' law (Eq.~\eqref{eq:stokes}), treating the shape factor $b_1$ as a free parameter. For the spheres, the best-fit value $b_1 = 17.7$ is consistent with the theoretical value of $18$ for smooth spheres (Fig.~\ref{fig:4}b), which lies within the confidence interval of the fit. This agreement indicates that the experimental setup accurately reproduces Stokes' settling behavior for smooth spheres. In contrast, silica grains settled more slowly, yielding $b_1 = 29.2$, a drag increase attributable to their angular, non-spherical shapes; \citeA{ferguson2004simple} reported a comparable $b_1 \approx 24$ for natural sediments of similar morphology.

Compared with silica, flocs settled more slowly and showed substantially greater dispersion in velocity at a given diameter (Fig.~\ref{fig:4}b). This reduced settling reflects the lower effective density of flocs arising from their highly porous, irregular internal structure \cite{kranenburg1994fractal, winterwerp1998simple}. Consequently, the floc data are not described by the standard Stokes formulation (Eq.~\eqref{eq:stokes}) and instead require the floc-specific relation (Eq.~\eqref{eq:floc_settling}), which introduces the 3D fractal dimension $n_f$ and the primary particle diameter $d_p$ as additional parameters.

\begin{figure}[htbp]
    \centering
    \includegraphics[width=1\textwidth]{figures/4.pdf}
    \caption{
    (a) Images of a laboratory sphere (Exp.~1), silica grain (Exp.~2), and a representative floc (Exp.~5). The 30~$\mu$m scale bar applies to all. (b) Log--log plot of settling velocity ($w_s$) vs.\ diameter ($d_f$) for spheres (green circles, Exp.~1), silica (red circles, Exp.~2), and aggregated flocs (blue circles, Exps.~3--8). Sphere and silica fits use Stokes' law (Eq.~\eqref{eq:stokes}), yielding shape factors ($b_1$) of 17.7 and 29.2. The floc data fit the fractal settling model ($w_s \propto d_f^{n_f - 1}$) with a slope of 0.90, giving an apparent fractal dimension $n_f = 1.90$. (c) Sensitivity of the fractal settling model (Eq.~\eqref{eq:floc_settling}) to assumed primary particle diameter ($d_p$). The floc data are compared to theoretical predictions using $n_f = 1.9$ for a primary particle size ($d_p$) of $10^{-6}$~m (dashed line) and $10^{-5}$~m (solid red line). This highlights how the choice of $d_p$ controls the required model fit, a key uncertainty in the conventional approach.
    }
    \label{fig:4}
\end{figure}

However, directly calculating $n_f$ by fitting Eq.~\eqref{eq:floc_settling} to the floc data is problematic because the primary-particle diameter, $d_p$, is unknown. Although our experiments have known distributions of input particle sizes (Fig.~\ref{fig:2}b), we do not know which particle sizes are actually incorporated into any given floc. Even if we could measure the primary particle sizes directly, the fractal theory necessitates specifying only a single value for a characteristic primary particle size, and it is unclear how to define this value from a distribution. Moreover, because $n_f$ appears both as an exponent on $d_f$ and within the pre-factor of Eq.~\eqref{eq:floc_settling} (via $d_p^{-(n_f-3)}$), a conventional regression that fits the model to the data for a prescribed $d_p$ (a one-parameter fit) forces the inferred $n_f$ to mathematically compensate for the assumed pre-factor. Consequently, the assumed value for $d_p$ exerts a strong influence on the inferred fractal dimension $n_f$ (Fig.~\ref{fig:4}c): assuming $d_p = 1\times 10^{-6}$~m yields $n_f = 2.77$; $d_p = 1\times 10^{-5}$~m gives $n_f = 2.10$; and $d_p = 1\times 10^{-4}$~m produces $n_f = 3.00$.

To avoid explicitly solving for the unknown primary-particle size $d_p$, we use the logarithmic form of Eq.~\eqref{eq:floc_settling}: $\log(w_s) = (n_f - 1)\log(d_f) + C$. Here, the intercept $C$ lumps together $d_p$, the shape factor $b_1$, and other constants. Under the assumption that \(d_p\) and \(b_1\) are approximately constant across the fitted population, the slope in \(\log(w_s)\)--\(\log(d_f)\) space yields \(n_f = \text{slope} + 1\) without requiring prior knowledge of \(d_p\). For the full floc population, the fitted slope is \(0.90\), giving \(n_f = 1.90\) (Fig.~\ref{fig:4}b). However, if structural heterogeneity causes \(n_f\) to covary with size, or if effective \(d_p\) varies across the population, a bulk regression can exhibit Simpson-type data-aggregation bias. We therefore stratify flocs by their \textit{in situ}, image-based fractal dimension and fit separate regressions within each stratum; the resulting slopes provide subgroup-specific estimates of \(n_f\) and reveal systematic dependence on structure.

\subsection{Estimating Fractal Dimension of Flocs from 2D Projections}

Fractal theory relates the fractal dimension to the primary particle diameter, the number of primary particles, and the floc diameter (Eq.~\eqref{eq:fractaltheory}). In natural systems, however, neither $d_p$ nor $N$ is constant or directly measurable during experiments. As a result, numerous methods have been developed to estimate floc fractal dimensions from image-based properties. These approaches are inherently two-dimensional, yielding an apparent fractal dimension $n_f^{(2\mathrm{D})}$ that may differ from the hydrodynamically relevant three-dimensional value. Throughout this paper, we denote the 3D fractal dimension simply as $n_f$, and use $n_f^{(2\mathrm{D})}$ when referring specifically to the 2D image-based measure.

The perimeter--area method has previously been used to estimate the 2D fractal dimension of flocs~\cite{lee2004prediction, maggi2004method, tang2015reconstructing}. This method relies on the empirical scaling relation:
\begin{equation}
P \propto A^{n_{f,\,\text{perim}}^{\text{(2D)}} / 2}
\end{equation}
where $P$ is the measured perimeter and $A$ is the projected area of the particle in the 2D image. In practice, $n_{f,\,\text{perim}}^{\text{(2D)}}$ is estimated from the log--log slope of $P$ versus $A$. This expression is derived by analogy to the Euclidean relation $P = k A^{1/2}$, which holds for smooth (non-fractal) boundaries with $n_f^{(2\mathrm{D})} = 1$. The underlying assumption is that, for a fractal boundary, the exponent in this scaling law approximates the true boundary fractal dimension $n_f^{(2\mathrm{D})}$. However, both mathematical analysis and geometric constructions (e.g., the Koch curve) demonstrate that this assumption breaks down for truly fractal sets~\cite{cheng1995perimeter,frame_mandelbrot_neger_2025}. For such boundaries, the measured perimeter is strongly dependent on the image resolution (the `yardstick' length), making a simple $P$--$A$ scaling ill-defined without explicitly accounting for the measurement scale. As a result, perimeter--area scaling does not reliably recover the true fractal dimension from 2D projections.

Box-counting, in contrast, is an appropriate method for fractal-dimension estimation that does not rely on strict self-similarity and is computationally straightforward \cite{foroutan1999advances}. A key advantage of box-counting is that it does not require knowledge of the primary particle size ($d_p$) or the number of constituent particles ($N$), making it particularly well-suited for natural flocs where these values are unknown or variable~\cite{bellouti1997flocs, spencer2022quantification}. The box-counting dimension $n_{f,\text{box}}^{\text{(2D)}}$ is determined by measuring how the number of occupied boxes $N_{\text{box}}(\epsilon)$ at a given scale $\epsilon$ varies with scale~\cite{mandelbrot1967long}:
\begin{equation}
N_{\text{box}}(\epsilon) \propto \epsilon^{-n_{f,\text{box}}^{\text{(2D)}}}
\end{equation}

For each floc mask, we computed $N_{\text{box}}(\epsilon)$ using square boxes with side lengths $\epsilon$ spanning powers of two in pixel units (\(\epsilon=1,2,4,\dots\)), and we estimated the dimension from the slope of a least-squares fit over the portion of the log--log curve that remains approximately linear and spans at least one order of magnitude in scale. When analyzing images of objects with finite extent from \textit{in situ} imaging, finite-size effects and image pixelation make the estimated fractal dimension sensitive to the available scaling range. For a digital image of $W \times W$ pixels, the practical range of box sizes ($\epsilon$) spans from 1 to $W$ pixels. Reliable estimation of $n_{f,\text{box}}^{\text{(2D)}}$ requires a sufficiently broad range of box sizes to establish a stable linear trend on a log--log plot of $N_{\text{box}}(\epsilon)$ versus $1/\epsilon$. If the image is too small, the limited scaling range can yield unstable values of $n_{f,\text{box}}^{\text{(2D)}}$. For example, a $20 \times 20$ pixel image provides only about 1.3 orders of magnitude in scale, which is insufficient for robust fitting and produces scale-dependent estimates (Fig.~\ref{fig:5}a). In contrast, a $100 \times 100$ pixel image spans two full orders of magnitude and typically provides 6--7 usable data points for regression; beyond this image size, the estimated $n_{f,\text{box}}^{\text{(2D)}}$ saturates, indicating a practical threshold for resolution-independent estimation (Fig.~\ref{fig:5}a). Accordingly, all analyzed flocs were extracted using $100 \times 100$ pixel windows centered on each particle. Fig.~\ref{fig:5}b shows examples for two experimental flocs at high image resolution, illustrating the application of the box-counting method to real data. In this study, all reported values of the 2D fractal dimension ($n_{f}^{\text{(2D)}}$) in the results and discussion sections refer exclusively to measurements obtained by the box-counting method.

\begin{figure}[htbp]
    \centering
    \includegraphics[width=1\textwidth]{figures/5.pdf}
    \caption{
    (a) Illustration of resolution dependence in box-counting dimension estimation. Top: A perfect circle embedded in $20 \times 20$ and $200 \times 200$ pixel grids shows that coarse grids systematically underestimate the fractal dimension. Bottom: Plot of measured 2D fractal dimension $n_f^{(2\mathrm{D})}$ versus image width $W$, highlighting a resolution-dependent regime at low $W$ and a stable region at high $W$. The transition scale ($W \approx 100$ pixels) was identified by repeating the box-counting calculation across increasing image sizes and selecting the smallest $W$ beyond which further increases produced negligible changes in the estimated dimension.
    (b) Example box-counting dimension estimates for two experimental flocs imaged at $>100 \times 100$ pixel resolution.
    }
    \label{fig:5}
\end{figure}

\subsection{Estimating Primary Particle Diameter ($d_p$) from Settling Data}

In most \textit{in~situ} experiments, the primary particle diameter \(d_p\) is assumed from the input sediment-size distribution and hypothesized floc-formation conditions because \(d_p\) is rarely resolvable directly. This practice can be uncertain when the effective primary scale differs from the feed distribution or varies across aggregation pathways. Here we identify \(d_p\) empirically by leveraging multiple floc subpopulations with distinct fractal dimensions \(n_f\). In the experimental application, we treat $d_p$ as an experiment-level parameter common to all $n_f$-stratified subpopulations. As implied by Eq.~\eqref{eq:floc_settling}, \(d_p\) does not affect the slope of the \(\log(w_s)\)--\(\log(d_f)\) relation (the slope is \(n_f-1\)), but it fixes a common intersection of the subgroup regressions. This is because at \(d_f=d_p\), the object is a primary particle (not a floc), so substituting \(d_f=d_p\) into Eq.~\eqref{eq:floc_settling} yields \(w_s=(g\,R_p/b_1\nu)\,d_p^{2}\), which is independent of \(n_f\). Consequently, regressions from any number of \(n_f\)-stratified subpopulations intersect at \((d_f, w_s)=(d_p,\,(g\,R_p/b_1\nu)\,d_p^{2})\). We refer to this common intersection diameter as $d_\text{intersect}$. In our synthetic example (Fig.~\ref{fig:1}), three subgroups with \(n_f \approx 1.5,\,2.0,\) and \(2.5\) produce regression lines that converge at \(d_\text{intersect} \approx 10^{-5}\,\mathrm{m}\), because the synthetic data were generated using this prescribed primary-particle diameter across varying fractal dimensions, thereby illustrating the method.


For typical observed flocs, the measured sizes imply that each floc contains many primary particles. For example, with $d_f/d_p \sim 10$--$100$ and $n_f \approx 2$, Eq.~\eqref{eq:fractaltheory} gives $N \sim 10^{2}$--$10^{4}$. This means the settling measurements reflect the collective behavior of aggregates built from a common sediment pool, and the inferred $d_p$ should be interpreted as an effective primary scale for the flocs resolved by the imaging and tracking. In the synthetic example, the subgroup regressions intersect exactly at a single point because the data are generated under the idealized assumption that all primary particles share a single diameter $d_p$. In natural sediments, primary particles span a distribution of sizes, so different flocs can be built from different effective primary scales and the subgroup regressions need not converge to a unique intersection. We therefore treat the intersection as distributed rather than singular and use a Monte Carlo procedure to generate an ensemble of intersection points, which we interpret as a distribution of effective $d_p$ values. We then compare this inferred $d_p$ distribution to the input sediment-size distribution to assess consistency and to quantify how the sediment's primary-size variability propagates into the settling-based estimate of $d_p$.

The intersection method assumes a common shape factor $b_1$ across fractal subgroups. Because drag depends on aggregate structure, $b_1$ may vary with $n_f$, biasing the inferred intersection diameter. Allowing for subgroup-dependent drag, the relationship between the observed intersection diameter $d_\text{intersect}$ and the true primary particle diameter $d_p$ for two subpopulations $j$ and $k$ is
\begin{equation}
d_\text{intersect} = d_p \left( \frac{b_{1,j}}{b_{1,k}} \right)^{1/(n_{f,j} - n_{f,k})},
\label{eq:dpintersection}
\end{equation}
which reduces to
\begin{equation}
d_\text{intersect} = d_p
\label{eq:dpintersectionideal}
\end{equation}
when $b_{1,j}=b_{1,k}$. In this limit, all subgroup regressions intersect exactly at the primary particle diameter.

\section{Results}

\subsection{Bulk Floc Size and Settling Velocity}

Before stratifying by fractal structure, we first report the bulk distributions of floc diameter ($d_f$) and settling velocity ($w_s$) across all experimental conditions (Table~\ref{tab:bulk_trends}). These summary statistics characterize the first-order response of the floc population to the imposed shear and organic matter forcing.

\begin{table}[htbp]
\caption{Median floc diameter and settling velocity with interquartile ranges (IQR) for each experimental condition.\label{tab:bulk_trends}}
\centering
\begin{tabular}{lcccccc}
\hline
Experiment & $u_*$ (m/s) & OM (wt\%) & $N$ & $d_f$ (\textmu m) [IQR] & $w_s$ (mm/s) [IQR] \\
\hline
LOWSHEAR  & 0.02 & 2 & 2429 & 43.2 [38.1--57.9] & 0.30 [0.21--0.44] \\
MIDSHEAR  & 0.03 & 2 & 2201 & 28.0 [24.6--35.9] & 0.39 [0.28--0.55] \\
HIGHSHEAR & 0.04 & 2 & 4603 & 26.8 [23.1--32.4] & 0.32 [0.24--0.46] \\
OM1       & 0.04 & 1 & 1058 & 42.0 [35.7--50.7] & 0.78 [0.63--1.00] \\
OM2       & 0.04 & 2 & 943  & 44.6 [36.8--55.6] & 0.67 [0.52--0.88] \\
OM3       & 0.04 & 3 & 1453 & 38.8 [32.0--50.5] & 0.52 [0.39--0.73] \\
\hline
\end{tabular}
\end{table}

In the shear-velocity series (Experiments 3--5; constant OM = 2~wt\%), median floc diameter decreases from 43.2~\textmu m at $u_* = 0.02$~m\,s$^{-1}$ to 26.8~\textmu m at $u_* = 0.04$~m\,s$^{-1}$, consistent with shear-driven breakup limiting maximum floc size \cite{winterwerp2002flocculation}. Despite the decrease in size, median settling velocity does not decrease proportionally: it is 0.30~mm\,s$^{-1}$ at low shear, rises to 0.39~mm\,s$^{-1}$ at mid shear, and returns to 0.32~mm\,s$^{-1}$ at high shear. This non-monotonic pattern reflects the competing effects of decreasing floc size and increasing floc compactness (higher $n_f$) with shear; smaller but denser flocs can maintain comparable settling velocities.

In the organic-matter series (Experiments 6--8; constant $u_* = 0.04$~m\,s$^{-1}$), median settling velocity decreases systematically from 0.78~mm\,s$^{-1}$ at 1~wt\% OM to 0.52~mm\,s$^{-1}$ at 3~wt\% OM, while median floc size also decreases modestly from 42.0 to 38.8~\textmu m. The decline in $w_s$ with increasing OM is consistent with flocs becoming more porous and less dense (lower $n_f$) as organic binding promotes open, voluminous structures. Thus, both size and structure change with OM, and both contribute to the observed settling-velocity trends.

These bulk trends motivate the stratified analysis that follows: the settling velocity depends not only on floc size but also on internal structure, and this structural variability must be resolved to accurately characterize $n_f$.

\subsection{Fractal Dimension from Stratified and Aggregated Analyses}

\subsubsection{Stratification Methodology and Bin Selection}

To infer a three-dimensional fractal dimension from settling data, we group flocs into bins based on their two-dimensional box-counting fractal dimension, $n_f^{(2\mathrm{D})}$, measured from images. The bin width in $n_f^{(2\mathrm{D})}$ is a coarse-graining scale rather than a physical control parameter. If a population contains discrete structural classes (e.g., non-flocculating grains with $n_f \approx 3$ and flocs with $n_f \approx 2$), finer stratification simply produces redundant bins that recover the same class-specific slopes. In finite datasets, the practical objective is therefore to resolve statistically distinct structural subpopulations while retaining enough particles per bin to constrain the conditional $w_s$--$d_f$ slope.

As an initial test case for flocs using our stratification methodology, we computed the two-dimensional box-counting fractal dimension, $n_f^{(2\mathrm{D})}$, for Experiment~3 (2~wt\% organic matter; shear velocity $u_* = 0.02~\mathrm{m\,s^{-1}}$). To evaluate the impact of stratification, we divided the dataset into bins of width 0.1, 0.05, and 0.025, corresponding to 3, 6, and 12 bins of $n_f^{(2\mathrm{D})}$, respectively, within a manually defined range from 1.3 to 1.6. Within each bin, we fitted a power-law relation between settling velocity, $w_s$, and floc diameter, $d_f$, in log--log space (Figs.~\ref{fig:6}a, d, g). The model $\log_{10}(w_s) = m \log_{10}(d_f) + C$ was calibrated using an ensemble Markov chain Monte Carlo (MCMC) sampler to obtain robust parameter estimates and verify convergence. After burn-in and thinning, the posterior medians of the slope $m$ were converted to the hydrodynamically relevant three-dimensional fractal dimension, denoted simply as $n_f$, using the relation $n_f = m + 1$.

Plotting the inferred three-dimensional fractal dimension, $n_f$, against the corresponding group-mean two-dimensional value, $n_f^{(2\mathrm{D})}$, shows how the stratification behaves across binning schemes. A linear trend captures the binned data reasonably well for all cases (Figs.~\ref{fig:6}c, f, i). However, because the number of bins is small (3, 6, or 12), the resulting coefficients of determination are not themselves a robust metric of model adequacy, as high $R^2$ values are expected even in the presence of substantial uncertainty.

A more informative assessment comes from quantifying the dynamical consequences of refining the stratification. From Eq.~\eqref{eq:floc_settling}, assuming the log--log regression lines converge near the primary particle diameter $d_p$, a change $\Delta n_f$ produces a vertical shift $\Delta \log_{10}(w_s) \approx \Delta n_f \log_{10}(d_f/d_p)$. Over the observed range of macroscopic floc sizes ($d_f/d_p \approx 10^1$--$10^2$), a change of $\Delta n_f = 0.01$ produces a shift of $\Delta \log_{10}(w_s) \approx 0.01$--$0.02$, corresponding to a change in settling velocity of less than 5\%. This effect is minimal relative to measurement scatter. Consequently, increasingly fine stratification does not reveal qualitatively different settling regimes but instead primarily inflates variance in the estimated slopes as bin populations decrease. We therefore adopt a bin width of 0.1 for subsequent analyses. This choice resolves meaningful structural heterogeneity while maintaining sufficient sample sizes for regression, and avoids over-resolving differences that have negligible impact on settling velocities over the observed floc-size range.

\begin{figure}[htbp]
    \centering
    \includegraphics[width=1\textwidth]{figures/6.pdf}
    \caption{Stratification of floc data by two-dimensional (2D) box-counting fractal dimension (\(n_f^{(2\mathrm{D})}\)) for Experiment~3, illustrating how binning strategy influences the estimation of three-dimensional fractal dimensions. (a--c) Analysis using 3 bins: (a) log--log plots of settling velocity (\(w_s\)) versus floc diameter (\(d_f\)), color-coded by 2D fractal-dimension bin; (b) fitted slopes (\(m\)) for each bin; (c) resulting three-dimensional fractal dimension (\(n_f = m + 1\)) plotted against mean \(n_f^{(2\mathrm{D})}\) for each bin, showing a strong linear relationship. (d--f) Same analysis using 6 bins; (g--i) using 12 bins, highlighting increased scatter and reduced regression stability with finer binning.}
    \label{fig:6}
\end{figure}

\subsubsection{3D Fractal Dimension Across All Experiments}

We expanded the analysis to all floc datasets (Experiments~3--8). Within each 0.1-wide \(n_f^{(2\mathrm{D})}\) bin, anomalous measurements were objectively identified and removed using an interquartile range (IQR) filter applied to the settling velocities: any particle with $\log_{10}(w_s)$ outside \([Q_1 - 1.5\,\mathrm{IQR},\ Q_3 + 1.5\,\mathrm{IQR}]\), where \(\mathrm{IQR} = Q_3 - Q_1\), was excluded. Only bins retaining at least 50 particles after filtering were analyzed, yielding a typical count of five to nine bins per experiment. Across experiments, analyzed bins typically contained tens to hundreds of flocs (median $\approx 280$ across bins with $n\ge50$), with peak counts concentrated near intermediate fractal dimensions (\(n_f^{(2\mathrm{D})}\approx1.4\!-\!1.6\)); bins near the distribution tails (\(n_f^{(2\mathrm{D})}\lesssim1.2\) or \(\gtrsim1.8\)) contained substantially fewer particles and were excluded by the minimum-count criterion. For each bin, the posterior median slope \(m\) from the log--log regression was converted to the corresponding three-dimensional fractal dimension as \(n_f = m + 1\). Inset panels in Fig.~\ref{fig:7} show the derived 3D \(n_f\) versus the bin-averaged \(n_f^{(2\mathrm{D})}\). Across all six experiments, stratified results reveal a strong positive correlation between these metrics, consistently supported by high \(R^2\) values.

The stratified analysis yields experiment-level median $n_f$ values that vary systematically with forcing. For the shear series (Fig.~\ref{fig:7}a--c), median $n_f$ increases from 1.77 at $u_* = 0.02$~m\,s$^{-1}$ to 1.95 at $u_* = 0.03$~m\,s$^{-1}$ and 2.12 at $u_* = 0.04$~m\,s$^{-1}$ (stratified IQRs of 0.26, 0.08, and 0.10, respectively). For the organic-matter series (Fig.~\ref{fig:7}d--f), median $n_f$ decreases from 1.84 at 1~wt\% OM to 1.66 at 2~wt\% OM and 1.63 at 3~wt\% OM (stratified IQRs of 0.14, 0.15, and 0.18). In contrast, the aggregated (bulk) analysis yields comparatively invariant median $n_f$ values: 1.98, 1.94, and 2.00 for the shear series, and 1.80, 1.74, and 1.75 for the OM series (all with narrow posterior IQRs of 0.03--0.04). The aggregated IQRs reflect only the posterior uncertainty of a single bulk $w_s$--$d_f$ regression, which concentrates with many observations; consequently, the aggregated approach assigns negligible probability to the lower $n_f$ values resolved by the stratified method.

\begin{figure}[htbp]
    \centering
    \includegraphics[width=1\textwidth]{figures/flocfigure7.pdf}
    \caption{Comparison of stratified and aggregated (bulk) fits for all floc datasets. (a--c) Varying shear velocities of 0.02, 0.03, and 0.04~m\,s$^{-1}$. (d--f) Varying organic matter contents of 1\%, 2\%, and 3\%. Inset plots show the relationship between \(n_f^{(2\mathrm{D})}\) and slope-inferred \(n_f\) for each experiment; points represent bin medians (with 95\% credible intervals) and are color-coded by \(n_f^{(2\mathrm{D})}\) group to make within-population variability visually explicit.}
    \label{fig:7}
\end{figure}

\subsubsection{Effects of Shear Velocity and Organic Matter}

The dependence of the stratified and aggregated $n_f$ estimates on environmental forcing is summarized in Fig.~\ref{fig:8}a,b. Stratified median $n_f$ increases approximately linearly with shear velocity at a rate of $\Delta n_f / \Delta u_* \approx 17$~s\,m$^{-1}$ ($n_f = 17.4\,u_* + 1.42$; Fig.~\ref{fig:8}a), consistent with shear-driven breakup preferentially destroying loose, low-$n_f$ flocs and favoring compact, high-$n_f$ survivors. The aggregated analysis yields a nearly flat response ($n_f = 1.0\,u_* + 1.94$), demonstrating that the shear-driven structural trend is effectively invisible in bulk data.

For organic matter (Fig.~\ref{fig:8}b), stratified median $n_f$ decreases at $\Delta n_f / \Delta \mathrm{OM} \approx -0.11$ per wt\% ($n_f = -0.11\,\mathrm{OM} + 1.92$), indicating that organic binding promotes open, porous structures. The aggregated trend is much weaker ($n_f = -0.03\,\mathrm{OM} + 1.81$). Stratified IQRs also widen with increasing OM (from 0.14 to 0.18), indicating that organic-rich populations contain enhanced structural diversity.

\begin{figure}[htbp]
    \centering
    \includegraphics[width=0.8\textwidth]{figures/flocpaperfigure8.pdf}
    \caption{Comparison of stratified and aggregated (bulk) analyses of floc fractal dimensions. (a) Median inferred \(n_f\) versus shear velocity. (b) Median \(n_f\) versus organic matter (OM) content. Boxes denote interquartile ranges (IQRs) and horizontal ticks indicate medians; for stratified estimates the IQR reflects particle-scale variability across bins weighted by population, whereas for aggregated estimates it reflects uncertainty in the Bayesian posterior of a single fitted parameter. (c) Relationship between three-dimensional (\(n_f\)) and two-dimensional (\(n_f^{(2\mathrm{D})}\)) box-counting fractal dimensions across all conditions, colored by experiment type. (d) Mapping between 2D and 3D fractal dimensions: blue dashed line, theoretical relation \(n_f = n_f^{(2\mathrm{D})} + 1\); red solid line, empirical 2D box-counting to 3D box-counting (BC) conversion ($n_f = 0.81\,(n_f^{(2\mathrm{D})})^{1.81}$); red dashed line, empirical 2D box-counting to 3D power-law (PL) conversion ($n_f = 0.20\,(n_f^{(2\mathrm{D})})^{4.08}$) \cite{wang2022fractal}. Experimental data (black markers) align most closely with the empirical box-counting model.}
    \label{fig:8}
\end{figure}

\subsubsection{Relationship Between 2D and 3D Fractal Dimensions}

Our experimental results (Fig.~\ref{fig:8}c) reveal a consistent positive relationship between the two- and three-dimensional fractal dimensions across all conditions. Plotting the stratified median \(n_f\) against the binned \(n_f^{(2\mathrm{D})}\) shows that denser, more compact structures (higher \(n_f^{(2\mathrm{D})}\)) systematically exhibit larger \(n_f\). This trend holds regardless of whether the grouping parameter is shear velocity or organic matter (OM) content.

To assess the underlying physical relationship, we compared our results to theoretical and empirical conversion models (Fig.~\ref{fig:8}d). The idealized linear relationship for perfectly space-filling Euclidean objects, \(n_f = n_f^{(2\mathrm{D})} + 1\) (blue dashed line), correctly describes smooth solids such as a sphere (\(n_f^{(2\mathrm{D})}=2\), \(n_f=3\)) but overpredicts the structural dimensions of porous flocs. The empirical models of \citeA{wang2022fractal}---a box-counting (BC) conversion, \(n_f = 0.81\,(n_f^{(2\mathrm{D})})^{1.81}\) (red solid line), and a power-law (PL) conversion, \(n_f = 0.20\,(n_f^{(2\mathrm{D})})^{4.08}\) (red dashed line)---provide better descriptions. Our data align most closely with the BC conversion, consistent with the sub-Euclidean space filling expected of porous aggregates: a 2D projection of a fractal object loses structural information, so the 3D dimension grows sublinearly rather than by a constant offset of unity. Although the data appear approximately linear within our observed range of \(n_f^{(2\mathrm{D})}\) (Fig.~\ref{fig:8}c), the underlying relationship is inherently nonlinear, as confirmed by the better agreement with the power-law BC conversion than with the linear Euclidean model.

\subsection{Inferring Primary Particle Diameter ($d_p$)}

To estimate the primary particle diameter ($d_p$), we used a Bayesian approach to quantify uncertainty in the intersection method. This approach assumes that the effective primary particle size controlling the settling relation is shared across the $n_f$-stratified subpopulations within a given run; accordingly, the inferred $d_p$ should be interpreted as an effective parameter in Eq.~\eqref{eq:floc_settling} rather than a direct census of the finest grains present in the source material. For each stratified floc group, we drew samples from the posterior distribution of the log--log settling fits and computed all pairwise intersections between fitted lines in log--log space. Under the constant-$b_1$ assumption (using $b_1 = 18$, the theoretical Stokes value for smooth spheres, as a reference baseline), the resulting intersection diameters $d_\text{intersect}$ provide an estimate of an effective $d_p$ (Eq.~\eqref{eq:dpintersectionideal}); allowing $b_1$ to vary among groups yields the more general relation in Eq.~\eqref{eq:dpintersection}. Repeating this calculation across all posterior draws yields an ensemble of intersection diameters; the resulting posterior distribution is what we plot as the inferred $d_p$ distribution in Fig.~\ref{fig:9}. All regression lines used in this intersection analysis satisfy the minimum sample-size criterion ($n\ge50$ flocs per group), and the posterior spread therefore reflects both measurement scatter and finite-sample uncertainty.

The intersection analysis yields inferred $d_p$ values of 21.7, 23.0, and 29.8~\textmu m for the low-, mid-, and high-shear experiments, and 34.2, 27.9, and 17.6~\textmu m for the 1\%, 2\%, and 3\% OM experiments.

\begin{figure}[htbp]
    \centering
    \includegraphics[width=1\textwidth]{figures/9.pdf}
    \caption{
        Cumulative percent finer by volume distributions comparing known input and inferred posterior effective primary particle diameters ($d_p$) for all experiments. Panels show (a--c) varying shear velocity and (d--f) varying OM content. The Kolmogorov--Smirnov (KS) statistic quantifies the agreement between the two distributions in each panel.
    }
    \label{fig:9}
\end{figure}

Fig.~\ref{fig:9} compares the cumulative percent finer by volume distributions for the known input particle sizes and the inferred $d_p$ from the intersection approach. The Kolmogorov--Smirnov (KS) statistic, reported in each panel, quantifies the maximum difference between the two cumulative distributions: a KS value of 0 indicates perfect agreement, while a value near 1 is a complete mismatch. Values between 0.29 and 0.41, as observed across all six experimental conditions, indicate moderate agreement, particularly in the central tendency (e.g., alignment of medians).

The input grain size distribution in Fig.~\ref{fig:9} shows all particle sizes added to the experiments, but not all of these necessarily serve as primary particles in flocs. Some input size fractions---such as the finest clays or coarsest silts and sands---may not be equally incorporated into flocs. Consequently, differences between the input grain size and inferred primary particle ($d_p$) distributions may reflect both uncertainties in the intersection method and real selective incorporation of certain sizes into flocs. Prolonged cloudiness after the experiments indicates that a fraction of very fine material remained in suspension; however, this observation alone cannot distinguish between incomplete aggregation of the smallest particles and limitations of the imaging and tracking workflow at small sizes. The inferred $d_p$ therefore represents an effective primary scale associated with the population of flocs contributing to the measured settling velocities, rather than the full input grain-size distribution. Because $d_p$ is obtained from regression-line intersections, an apparent truncation of the fine tail reflects a concentration of the mass contributing to the observed $w_s$--$d_f$ scaling within a restricted size range under the experimental freshwater chemistry and shear conditions. The offset between input and inferred distributions may arise from size-dependent participation in the floc population resolved by the measurements and from under-sampling of dispersed particles below the detection threshold.

The shape factor $b_1$ may vary with fractal structure, which would cause $d_\text{intersect}$ to deviate from the true $d_p$. To illustrate this effect, we plot settling velocity versus diameter for two fractal groups ($n_f = 1.5$ and $n_f = 2.0$), both with a common primary particle diameter of $10~\mu$m, using Eq.~\eqref{eq:floc_settling} (Fig.~\ref{fig:10}a). For $n_f = 1.5$, we fix $b_1 = 50$, and for $n_f = 2.0$, we vary $b_1$ from 10 to 100. As the shape factor increases with fractal dimension, the intersection diameter $d_\text{intersect}$ between the two groups becomes greater than the true primary particle diameter. The opposite holds if $b_1$ decreases with $n_f$.

\begin{figure}[htbp]
    \centering
    \includegraphics[width=1\textwidth]{figures/10.pdf}
    \caption{
    Influence of the shape factor (\(b_1\)) on the geometric intersection diameter (\(d_\text{intersect}\)) between floc groups of different fractal dimension.
    (a) Increasing \(b_1\) with \(n_f\) shifts \(d_\text{intersect}\) above the primary particle diameter \(d_p\), while decreasing \(b_1\) shifts it below.
    (b) When evaluating the idealized intersection assuming constant \(b_1\), the lower tail of measured floc diameters extends to the left of \(d_\text{intersect}\). This violates the physical constraint that flocs cannot be smaller than their fundamental building blocks.
    (c) Example where \(d_\text{intersect} > d_p\), emphasizing that \(b_1\) must rise with \(n_f\) to preserve the proper size hierarchy as denser flocs experience greater drag.
    }
    \label{fig:10}
\end{figure}

We evaluate the calculated geometric intersection diameters ($d_\text{intersect}$) obtained under the standard assumption of a constant shape factor (Fig.~\ref{fig:10}b). We find that while most measured floc diameters lie to the right of $d_\text{intersect}$, the lower tail of the observed floc size distribution inevitably extends to the left (i.e., some observed $d_f < d_\text{intersect}$). Because a floc physically cannot be smaller than its constituent primary particles ($d_f \ge d_p$), the true primary particle diameter must be strictly smaller than all measured floc diameters. Consequently, the assumption that $d_\text{intersect} = d_p$ produces a physical contradiction, requiring instead that the true primary particle size be securely anchored to the left of the constant-$b_1$ intersection ($d_p < d_\text{intersect}$).

This strict physical constraint implies a necessary relationship between $b_1$ and $n_f$. According to Eq.~\eqref{eq:dpintersection}, the only way to mathematically maintain $d_p < d_\text{intersect}$ as $n_f$ increases is if $b_1$ also increases. This requirement ensures that the primary particle remains the smallest logical constituent in the system. Such a trend is consistent with theoretical expectations and experimental observations, which suggest that denser, less permeable flocs (i.e., those with higher $n_f$) experience greater hydrodynamic drag, corresponding to higher values of $b_1$ \cite{strom2011explicit}. Importantly, this conclusion does not invalidate the $n_f$ values inferred from the stratified $\log(w_s)$--$\log(d_f)$ slopes, because those slopes are independent of $b_1$; the shape factor enters only the intercept. The sensitivity of $d_\text{intersect}$ to $b_1$ therefore affects only the inferred primary particle diameter, not the fractal dimension itself.

\section{Discussion}

\subsection{How Fractal Structure Mediates Settling Velocity}

A central finding of this study is that the 3D fractal dimension varies systematically with environmental forcing and that this variation has direct consequences for settling velocity. The stratified analysis reveals that $n_f$ increases with shear velocity (from 1.77 to 2.12 over $u_* = 0.02$--$0.04$~m\,s$^{-1}$) and decreases with organic matter content (from 1.84 to 1.63 over 1--3~wt\% OM). These trends reflect two competing physical processes that govern floc structure: shear-driven breakup, which preferentially destroys fragile, open aggregates and selects for compact survivors with higher $n_f$; and organic binding, in which extracellular polymeric substances act as a flexible ``glue'' that holds particles in extended, porous configurations with lower $n_f$ \cite{nghiem2022mechanistic}.

These structural changes propagate directly into settling velocity. From Eq.~\eqref{eq:floc_settling}, $w_s$ depends on floc size as $d_f^{n_f-1}$ and on floc density (via $d_p^{3-n_f}$), so both size and $n_f$ jointly determine the settling rate. In the shear series, increasing $u_*$ simultaneously decreases floc size (median $d_f$ drops from 43 to 27~\textmu m; Table~\ref{tab:bulk_trends}) and increases compactness ($n_f$ rises from 1.77 to 2.12). These effects partially offset each other: smaller flocs settle more slowly, but denser flocs settle faster. The net result is that median $w_s$ remains roughly comparable across shear levels ($0.30$--$0.39$~mm\,s$^{-1}$), consistent with the non-monotonic bulk trend in Table~\ref{tab:bulk_trends}. In the OM series, increasing organic content both reduces compactness ($n_f$ drops from 1.84 to 1.63) and modestly decreases floc size (median $d_f$ drops from 42 to 39~\textmu m). Both effects reduce $w_s$, and their reinforcement drives the pronounced decline from 0.78 to 0.52~mm\,s$^{-1}$.

In contrast, the aggregated (bulk) approach yields $n_f \approx 1.9$--$2.0$ across all conditions and would predict that changes in $w_s$ are driven almost entirely by size. The stratified analysis reveals that this interpretation is incomplete: structural changes account for a substantial fraction of the settling-velocity response, particularly in the OM series where the stratified $n_f$ drops well below the aggregated estimate.

These trends are broadly consistent with prior work. \citeA{winterwerp1998simple} reported $n_f$ as low as 1.4 for fragile marine-snow aggregates and up to 2.2 for robust estuarine flocs, and attributed the range to differences in shear history and biological binding. Our results extend this framework by quantifying, within a single experimental system, the rates at which $n_f$ responds to shear ($\approx$17~s\,m$^{-1}$) and organic matter ($\approx -0.11$ per wt\%), and by demonstrating that these responses are detectable only when structural heterogeneity is resolved.

\subsection{2D to 3D Conversion of Fractal Dimension}

The stratified data reveal a positive, nonlinear relationship between 2D box-counting fractal dimension and the hydrodynamically inferred 3D fractal dimension (Fig.~\ref{fig:8}c,d). The idealized Euclidean conversion, $n_f = n_f^{(2\mathrm{D})} + 1$, systematically overpredicts $n_f$ because porous aggregates do not fill space completely: a 2D projection of a fractal object loses structural information, so the 3D dimension grows sublinearly. Our data align most closely with the empirical box-counting conversion of \citeA{wang2022fractal} ($n_f = 0.81\,(n_f^{(2\mathrm{D})})^{1.81}$), derived from computationally generated aggregates.

Although this agreement is encouraging, no universally valid mapping between 2D and 3D fractal dimensions exists for natural flocs. Synthetic studies, including \citeA{wang2022fractal}, grow aggregates from monodisperse primary particles with fixed size, whereas natural flocs contain a broad distribution of particle sizes and shapes, violating the fixed-$d_p$ assumption. These differences introduce scatter and potential systematic offsets. Furthermore, permeability---which is not considered here---can alter settling without changing geometry, adding an additional unmapped degree of freedom. The 2D--3D relationship should therefore be treated as system specific and re-evaluated when sediment mineralogy, salinity, or aggregation pathways differ substantially. Nevertheless, the 2D box-counting dimension remains a practical stratification tool: by quantifying within-population variability and enabling direct experimental calibration, our approach provides context-sensitive inference of the hydrodynamically relevant 3D fractal dimension.

\subsection{Constraining Primary Particle Diameter ($d_p$) and Shape Factor ($b_1$)}

The primary particle diameter ($d_p$) and the shape factor ($b_1$) are key parameters in fractal settling theory, yet they are rarely measured directly. We estimate $d_p$ by identifying $d_\text{intersect}$, the diameter at which the $n_f$-stratified log--log regression lines converge (Fig.~\ref{fig:9}). In principle, at $d_f = d_p$ all floc populations should converge to the settling velocity of a single primary particle regardless of $n_f$.

The intersection analysis yields median $d_\text{intersect}$ values that agree with the median of the input sediment sizes (Fig.~\ref{fig:9}), yet the inferred distribution is narrower: about two orders of magnitude wide compared with nearly four orders in the input material. This discrepancy may arise because only a subset of the source material (coarse silt to fine sand) effectively participates in forming the flocs that dominate the measured settling velocities, or because the imaging system does not capture the smallest dispersed grains. \citeA{nghiem2024testing} showed that the effective $d_p$ in fractal scaling is better represented by a number-weighted moment of the size distribution rather than a volumetric median, because physical aggregation conserves particle number rather than volume. The inferred $d_p$ values therefore represent an effective primary scale for the resolved floc population rather than a census of all grains present.

A key result of this analysis is the demonstration that $b_1$ must increase with $n_f$. Under the assumption of constant $b_1$, the intersection diameter $d_\text{intersect}$ exceeds some observed floc diameters (Fig.~\ref{fig:10}b), which is physically impossible because flocs cannot be smaller than their building blocks. Eq.~\eqref{eq:dpintersection} shows that $d_\text{intersect} > d_p$ requires $b_1$ to increase with $n_f$: denser, less permeable flocs experience greater form drag and thus higher effective $b_1$ \cite{strom2011explicit}. This finding is consistent with theoretical expectations but, to our knowledge, has not previously been demonstrated empirically from settling data alone. We emphasize that this $b_1$--$n_f$ dependence affects only the inferred $d_p$ (via the regression intercept) and does not alter the slope-based $n_f$ estimates, which are independent of $b_1$.

Uncertainty in $d_p$ propagates directly into predictions: the value of $n_f$ inferred from a full-model fit to Eq.~\eqref{eq:floc_settling} is highly sensitive to the assumed $d_p$ (Fig.~\ref{fig:4}c). The slope-based approach adopted here avoids this coupling, but predictive use of Eq.~\eqref{eq:floc_settling} ultimately requires both parameters. In practice, the effective $d_p$ relevant for flocculation can be both narrower in range and offset in value relative to the bulk sediment size distribution.

\subsection{Implications for Sediment Transport Modeling}

The 3D fractal dimension $n_f$ governs the density--size relationship of flocs and thereby controls settling velocity in morphodynamic models. Historically, modelers have adopted $n_f \approx 2.0$--$2.2$ from bulk $w_s$--$d_f$ regressions \cite{kranenburg1994fractal, son2008flocculation, winterwerp1998simple, winterwerp2004introduction, winterwerp2006heuristic}. Our aggregated analysis yields $n_f \approx 2.0$, confirming that our experimental flocs are physically comparable to those in previous studies. However, the stratified analysis reveals that this apparent consistency masks substantial structural diversity.

By resolving this diversity, the stratified method yields median $n_f$ values of 1.6--2.1 depending on forcing, with the majority of conditions producing $n_f < 2.0$. These lower fractal dimensions correspond to more porous, less dense flocs that settle more slowly. For a representative 400~\textmu m floc with 4~\textmu m primary particles, lowering $n_f$ from 2.2 to 1.7 in Eq.~\eqref{eq:floc_settling} produces an order-of-magnitude decrease in settling velocity, and a corresponding order-of-magnitude increase in transport distance. In many transport formulations, the deposition flux scales as $F_{\mathrm{dep}} \sim w_s\,C$ (with suspended concentration $C$), so lowering $n_f$ and thus $w_s$ reduces modeled deposition rates and increases downstream export for the same supply \cite{amoudry2008review}. The efficiency of mud retention in freshwater floodplains and deltas may therefore be considerably lower than models assuming $n_f \geq 2.0$ would predict.

For practitioners, we recommend adopting $n_f$ in the range of 1.6--2.0, and, where feasible, representing the fractal dimension as a distribution rather than a single value. Selecting a lower $n_f$ yields predictions of reduced deposition and increased sediment bypass relative to the conventional range. While next-generation models should ideally represent the full distribution of $n_f$ to capture natural heterogeneity, a pragmatic immediate step is to shift default parameterizations away from the biased, artificially high values that arise from bulk data aggregation. This adjustment is important for improving predictions of fine-sediment fate in freshwater systems; extending the framework to saline estuaries requires explicit consideration of salt-driven physical aggregation, which may shift the $n_f$ distribution toward higher values.

\section{Conclusions}

This study demonstrates that accounting for structural heterogeneity within floc populations is essential for accurate sediment transport modeling. Aggregating settling data across flocs with varying fractal dimensions produces misleading results due to Simpson's Paradox: physical trends present within distinct structural subgroups are lost or distorted when data are analyzed in bulk. By stratifying flocs based on 2D box-counting fractal dimensions, we resolve these structural differences and enable physically meaningful inference of floc transport parameters.

The stratified approach reveals that the hydrodynamically inferred 3D fractal dimension ($n_f$) increases with shear velocity ($u_*$) and decreases with organic matter content, reflecting the competition between shear-driven breakup---which selects for compact, high-$n_f$ structures---and organic binding, which promotes open, porous aggregates with low $n_f$. These trends, quantified as $\Delta n_f / \Delta u_* \approx 17$~s\,m$^{-1}$ and $\Delta n_f / \Delta \mathrm{OM} \approx -0.11$ per wt\%, are invisible in bulk analyses that yield a near-constant $n_f \approx 2.0$. The relationship between 2D and 3D fractal dimensions is inherently nonlinear, aligning with the empirical box-counting conversion of \citeA{wang2022fractal} rather than the idealized Euclidean offset, consistent with the sub-Euclidean space filling of porous aggregates. Furthermore, the stratified intersection analysis constrains the effective primary particle diameter and demonstrates that the hydrodynamic shape factor ($b_1$) must increase with $n_f$ to maintain physical consistency.

These results show that the widespread use of $n_f \approx 2.0$--$2.2$ in sediment transport models overestimates floc compactness, leading to overprediction of settling velocity by up to an order of magnitude and corresponding underprediction of transport distance. We recommend adopting $n_f$ in the range 1.6--2.0, which better represents the porous, slow-settling flocs that dominate suspended transport but remain hidden in bulk analyses. Making this adjustment will improve predictions of fine-sediment transport, deltaic development, and downstream sediment budgets in freshwater-dominated systems.

\section*{Open Research Section}
All code and analysis workflows supporting the results and conclusions of this study are openly available at \url{https://github.com/braydennoh/FlocTrack}. The repository includes scripts for image analysis, fractal dimension calculation, and all processing steps described in this manuscript. For further questions or access to additional datasets, please contact the corresponding author at braydennoh@fas.harvard.edu.

\section*{Inclusion in Global Research Statement}
This research did not involve collaborations with local communities or cross-regional partnerships requiring additional disclosure. All research activities were conducted in accordance with institutional and international ethical standards.

\acknowledgments
BN acknowledges support from Caltech SURF, the Arthur R. Adams SURF Fellowship, and the Saul and Joan Cogen Memorial SURF Fellowship at Caltech. MPL acknowledges support from NASA FINESST Grant 80NSSC20K1645, the NASA Delta-X project funded by the Science Mission Directorate's Earth Science Division through the Earth Venture Suborbital-3 Program NNH17ZDA001N-EVS3, and the National Science Foundation Geomorphology and Land-use Dynamics Grant No. 2136991.

\bibliography{agusample}


\end{document}
